\documentclass[]{article}

\usepackage{graphicx}
%opening
\title{\textbf{RESUME DES EXPOSES D'INVESTIGATION NUMERIQUE}}

\date{}
\begin{document}
\thispagestyle{empty}
\begin{center}
\rule{1\textwidth}{2pt}\vspace{\baselineskip}
\vspace{\baselineskip}\huge\textbf{THEORIES ET PRATIQUES DE L’INVESTIGATION  NUMERIQUE.}
\rule{1\textwidth}{2pt}\vspace{\baselineskip}
\vspace{\baselineskip}
\vfill
\textbf{Enseignant: M.MINKA THIERRY.}\\
\vfill
Par l'étudiante: CHEUTCHEULONGA HILARY PAULA\\
Filière CIN 4\vfill
\vspace{\baselineskip}
\emph{\textbf{Année Academique: 2025-2026.}}
\rule{1\textwidth}{2pt}\vspace{\baselineskip}
\end{center}
\maketitle
\section{DEEPFAKE VIDEO}
\paragraph{}
L'émergence de l'intelligence artificielle générative a profondément transformé les pratiques de création de contenus numériques, en offrant de nouvelles possibilités dans des domaines aussi variés que la communication, la pédagogie, le divertissement ou encore la recherche. Parmi ces avancées, l'association des modèles de langage, capables de produire des scripts structurés et cohérents, et des plateformes de synthèse vidéo, spécialisées dans la génération de visuels réalistes, constitue un axe d'expérimentation particulièrement prometteur. C'est dans ce contexte que s'inscrit le travail du premier groupe : la réalisation d'un deepfake pédagogique, reposant sur l'utilisation de GPT pour la rédaction du script et de Heygen pour la génération vidéo. Ce projet illustre concrètement comment l'IA peut être mise au service de la production de contenus immersifs, tout en soulevant des questions éthiques et méthodologiques.Nous apprenons qu’un un deepfake est un enregistrement vidéo ou audio réalisé ou modi é grâce à l'intelligence artificielle. Ce terme fait référence non seulement au contenu ainsi créé, mais aussi aux technologies utilisées. Le mot deepfake est une abréviation de "Deep Learning" et "Fake", qui peut être traduit par "fausse profondeur". En fait, il fait référence à des contenus faux qui sont rendus profondément crédibles par l'intelligence artificielle.\\ En 2014, une technique inventée par le chercheur Ian Goodfellow est à l'origine des deepfakes. Il s'agit du GAN (Generative Adverarial Networks). Selon cette technologie, deux algorithmes s'entraînent mutuellement : l'un tente de fabriquer des contrefaçons aussi fiables que possible ; l'autre tente de détecter les faux. De cette façon, les deux algorithmes s'améliorent ensemble au fil du temps grâce à leur entraînement respectif. Plus le nombre d'échantillons disponibles augmente, plus l'amélioration de ceux-ci est importante. Pour réaliser cette vidéo, des outils telsque HEYGEN AI peut etre utilisé pour faire une fausse video. HeyGen est une intelligence artificielle spécialisée dans la génération de vidéos par IA. Créé en 2022, cet outil se distingue par sa capacité à transformer de simples instructions textuelles en séquences audiovisuelles captivantes, sans exiger de matériel sophistiqué ni de compétences techniques avancées. Son accessibilité et sa rapidité d'exécution ouvrent un large éventail de possibilités, aussi bien pour les professionnels que pour les particuliers, en permettant de produire en quelques clics des contenus de qualité quasi professionnelle.

\section{SIMILATION D’UNE SERIE DE MESSAGE SUR WHATSAPP ENTRE UN HOMME ET SA MAITRESSE}
\paragraph{}
Dans le contexte actuel où les échanges numériques occupent une place centrale dans
la vie sociale et personnelle, les applications de messagerie instantanée comme WhatsApp
constituent une source d'information privilégiée, mais aussi un vecteur de manipulations.
L'investigation numérique vise précisément à analyser, comprendre et parfois reproduire
ces environnements afin d'étudier les traces laissées par les utilisateurs ou de mettre en
évidence la possibilité de falsifications. Dans le cadre de ce travail, nous avons choisi de
simuler une série de messages échangés entre un homme et sa maîtresse, en mobilisant
deux outils : Chatsmock, qui permet de générer de fausses conversations WhatsApp, et
Adobe Photoshop, utilisé pour affiner et personnaliser l'apparence des échanges. L'objectif
n'est pas de porter un jugement moral sur les faits simulés, mais d'illustrer les possibilités
techniques offertes par ces logiciels et de questionner la fiabilité des preuves numériques
dans un contexte d'enquête. Le scénario imaginé est celui d’un enseignant qui entretiendrait une relation extra-conjugale avec une étudiante de son établissement, afin de mieux comprendre les enjeux liés à la manipulation des preuves numériques, une simulation de conversations a été réalisée sur l'application WhatsApp. Chatsmock est une application web permettant de créer de fausses conversations WhatsApp de manière réaliste. Grâce à son interface intuitive, il est possible de :\\
- Définir les participants à la discussion (noms, photos de pro l, numéros de télé-
phone).\\
- Générer des messages avec un contenu librement choisi.\\
- Paramétrer l'heure, la date et le statut de lecture (message envoyé, reçu, ou vu).\\
- Créer des captures d'écran identiques à celles produites par l'application WhatsApp
réelle.\\
Afin d'améliorer le réalisme des captures produites, Adobe Photoshop a ensuite été
utilisé pour : Corriger certains détails graphiques (alignement, bulles de texte, couleur des icônes). Modifier ou insérer des éléments supplémentaires, comme des images envoyées dans
la discussion. Retoucher l'interface visuelle a n qu'elle corresponde parfaitement à celle d'un
smartphone réel.

\section{POINTS SUR L’ALGORITHME DE RECONNAISSANCE FACIALE}
\paragraph{}
La reconnaissance faciale est une technologie d'intelligence artificielle permettant d'identifier ou de vérifier l'identité d'une personne à partir de ses traits du visage. Elle repose sur des algorithmes capables de détecter, d'analyser et de comparer des caractéristiques faciales uniques (par exemple la distance entre les yeux, la forme du nez, les contours de la mâchoire ou des lèvres). Ces algorithmes sont largement utilisés dans divers domaines : sécurité (vidéosurveillance, contrôle d'accès), téléphonie mobile (déverrouillage de smartphones), réseaux sociaux (étiquetage automatique de photos), etc. Toutefois, leur utilisation soulève des enjeux éthiques et juridiques importants (protection des données personnelles, respect de la vie privée). La reconnaissance faciale (RF) est une technique biométrique qui utilise le visage d'une personne comme trait biométrique. Elle est très répandue du fait que le visage est une modalité biométrique non intrusive et facile à acquérir. Un système biométrique fonctionne en trois phases principales : l'enrôlement, l'identification et la vérification.\\ 
-	Lors de l'enrôlement, l'utilisateur est inscrit pour la première fois : sa modalité biométrique (le visage) est capturée, prétraitée, et les caractéristiques pertinentes en sont extraites puis stockées dans une base de données. Des informations biographiques peuvent être associées à ce profil.\\ 
-	En mode identification, un individu se présente au système sans révéler explicitement son identité ; le système cherche alors parmi tous les profils enrôlés le plus proche par correspondance (recherche1-N).\\
-	En vérification (authentification), l'utilisateur affirme une identité, et le système compare les caractéristiques extraites à son modèle enregistré (recherche 1-1). Si le score de similitude est supérieur à un seuil, l'identité est confirmée.\\

\section{DEEPFAKE VOCAL}
\paragraph{}
Un deepfake en français faux profond selon Fortinet est une forme d’intelligence artificielle qui peut être utilisée pour créer des images, sons et des vidéos de canulars convaincants. Dans ce cas, nous travaillerons sur un deepfake audio, c’est-à-dire un enregistrement sonore falsifié produit à l’aide de techniques d’apprentissage profond (deep learning). Il consiste à entraîner un modèle d’intelligence artificielle sur un ensemble de voix réelles pour ensuite imiter la voix d’une personne cible et générer des paroles qu’elle n’a jamais prononcées. Les deepfakes audio illustrent parfaitement le défi contemporain entre progrès technologique et préservation de l’intégrité des preuves. En menaçant la confidentialité, la fiabilité et l’opposabilité, ils imposent aux experts de renforcer leurs protocoles d’authentification et de maintenir une transparence irréprochable. Comprendre leur fonctionnement n’est plus un atout : c’est une condition essentielle pour protéger la crédibilité des preuves et préserver la confiance dans l’investigation numérique.\\
La technologie utilisé dans ce cas est celui de MINIMAX audio, qui est une technologie basée sur l’intelligence artificielle et l’apprentissage profond, spécialisée dans la synthèse et la transformation de la voix humaine. Elle utilise des modèles entraînés sur de vastes bases de données vocales pour imiter avec précision le timbre, l’intonation et le rythme d’un locuteur réel. L’objectif est de proposer des applications innovantes dans les domaines du divertissement, de l’éducation et de l’accessibilité numérique. Par exemple, elle permet de recréer des voix pour des doublages de films, des podcasts ou encore pour des assistants vocaux personnalisés.
Cependant, les deepfakes vocaux suscite aujourd’hui de vives inquiétudes dans les domaines
de la cybersécurité, de la justice et des médias. Plusieurs cas réels illustrent déjà les dangers de cette technologie.\\
•	En 2019, un dirigeant d’une entreprise britannique d’énergie a été trompé par un deepfake vocal imitant la voix de son PDG allemand, l’amenant à transférer plus de 220 000 euros sur un compte frauduleux (Forbes, 2019).\\
•	En 2020, une étude menée par l’université de Stanford a montré que des voix clonées pouvaient tromper jusqu’à 50 pourcent des systèmes de reconnaissance vocale commerciaux (Stanford HAI, 2020).\\
•	Plus récemment, en 2022, des chercheurs en cybersécurité ont démontré qu’il était possible de cloner une voix à partir de seulement 5 secondes d’enregistrement, ce qui ouvre la voie.\\
Face aux menaces posées par le clonage vocal, plusieurs solutions émergent et doivent être
appliquées :Développement d’outils capables d’analyser les signaux vocaux pour identifier des anomalies propres aux voix générées par IA. Ces détecteurs pourraient être intégrés dans les plateformes de communication tels que les réseaux sociaux. Les utilisateurs doivent être formés pour reconnaître les risques. Dans les entreprises, des programmes de prévention peuvent réduire les fraudes basées sur les deepfakes vocaux. Pour le cadre légal et réglementaire, plusieurs pays réfléchissent à des lois spécifiques sur les deepfakes, imposant des sanctions et un marquage numérique (watermarking) des contenus générés. La mise en place de méthodes d’authentification robustes : reconnaissance vocale dynamique et combinaison avec l’authentification multi-facteur.

\section{CONCEPTION ET ANALYSE D’UN FAUX PROFIL TIKTOK : CHOIX D’UNE NICHE DANS LE CADRE D’UNE INVESTIGATION NUMÉRIQUE.}
\paragraph{}
Les réseaux sociaux façonnent l’opinion, influencent les comportements et redéfinissent les interactions humaines, TikTok s’impose comme une plateforme incontournable, notamment auprès des jeunes générations. Notre travail d’investigation numérique s’inscrit dans une démarche pédagogique visant à comprendre les enjeux liés à l’identité numérique, à la viralité des contenus et aux risques de manipulation de l’information. Dans ce cadre, un faux profil a été créé sur TikTok autour de la thématique de la cybersécurité, une niche à la fois technique, sensible et d’actualité. L’objectif étant d’analyser les réactions suscitées, les types d’interactions générées, mais aussi de réfléchir aux implications éthiques d’un tel projet.\\
Dans le cadre notre travail, un faux profil tiktok a été crée pour débuter. Cette étape a nécessité l’utilisation d’un service de messagerie temporaire (TEMP MAIL) afin de préserver l’anonymat et d’éviter toute liaison directe avec nos données personnelles. Ce profil nous a servi de base pour observer et analyser les interactions au sein de la niche choisie la cybersécurité. Pour assurer un suivi efficace de l’activité du faux profil TikTok, plusieurs outils numériques ont été mobilisés, dans une optique d’observation et d’analyse du comportement des utilisateurs face à nos contenus :Utilisation des statistiques internes de TikTok (TikTok Analytics) : nombre de vues, de likes, de partages, taux d’engagement, évolution du nombre d’abonnés, etc. ; Captures d’écran effectuées à différentes étapes de la vie du profil ; Générateurs de contenu comme ChatGPT pour les messages, Canva pour les visuels, Temp Mail pour le compte ; WhatsApp pour le partage du lien d’accès au compte ; Un tableau de bord personnel pour noter observations et hypothèses. Ce projet met en lumière le pouvoir de manipulation des réseaux sociaux, tout en rappelant la responsabilité qui accompagne la diffusion de contenus, même fictifs. L’usage de fausses identités, même pour des fins de sensibilisation, doit toujours s’encadrer strictement dans une démarche responsable et éthique.

\section{LES TROIS MEILLEURS LOGICIELS DE RÉDACTION DE MÉMOIRE}
\paragraph{}
La rédaction d'un mémoire représente un défis académique majeur, tant par son ampleur que par sa complexité. Entre la gestion fastidieuse des sources bibliographiques, le respect des normes formelles et la structuration d'un contenu substantiel, l'étudiant se trouve confronté à une entreprise qui dépasse largement le cadre de la simple rédaction. Dans ce contexte exigeant, le choix des outils logiciels devient un paramètre déterminant pour la réussite du projet. Un "bon" logiciel pour mémoire doit répondre à plusieurs impératifs : offirir un environnement de rédaction adapté aux longs documents, que ce soit via LaTeX ou un traitement de texte classique ; garantir une gestion rigoureuse des références bibliographiques ; et faciliter la mise en forme selon les standards académiques. Face à la multitude d'outils disponibles, se pose alors cette question cruciale : quels sont les logiciels les plus efficaces et les mieux adaptés pour accompagner l'étudiant dans la réussite de son mémoire ? Pour y répondre, trois types d'outils spécialisés sont analyse successivement : un environnement LaTeX professionnel avec Overleaf, une solution de traitement de texte universelle avec Word, et un gestionnaire de références bibliographiques avec Zotero.\\
- Overleaf se définit comme un éditeur LaTeX en ligne collaboratif qui a révolutionné l'approche de la rédaction académique. Sa philosophie repose sur trois piliers fondamentaux : L'accessibilité qui rendre LaTeX utilisable sans installation complexe ; La collaboration qui permettre un travail d'équipe fluide et synchronisé ; La qualité : maintenir les standards professionnels de l'édition scientifique. Particulièrement standard dans les disciplines scientifiques et techniques (mathématiques, physique, informatique, ingénierie), il incarne l'approche "what you mean is what you get" - où l'auteur se concentre sur le contenu sémantique tandis que le système gère automatiquement la mise en forme professionnelle.\\
- Microsoft Word s'est imposé comme l'outil de rédaction universel, détenant une position quasi hégémonique dans le monde académique et professionnel. Sa prédominance s'explique par plusieurs facteurs : son intégration à la suite Microsoft Office, présente dans la plupart des institutions éducatives ; son interface graphique familière pour la majorité des utilisateurs, qui l'ont généralement découvert dès leur scolarité ; et sa compatibilité quasi universelle, faisant du format .docx un standard d'échange incontournable. Cette familiarité en fait souvent le choix par défaut, voire le seul logiciel de traitement de texte maîtrisé par de nombreux étudiants et encadrants.\\
- Zotero se présente comme un gestionnaire de références bibliographiques open-source
spécialement conçu pour la recherche académique. Son approche repose sur une philosophie de centralisation des références : offrir un écosystème unique où l'étudiant peut rassembler, organiser et exploiter l'ensemble des sources de son mémoire. Développé initialement par le Center for History and New Media de l'Université George Mason, Zotero incarne une vision démocratique de la gestion bibliographique, combinant puissance fonctionnelle et gratuité. Son modèle open-source lui permet une grande flexibilité et une communauté active de contributeurs.\\
Chaque solution présente des atouts complémentaires : Overleaf excelle par sa qualité typographique irréprochable et son approche structurante, idéale pour les documents complexes. Microsoft Word conserve son avantage majeur en termes d'accessibilité et de prise en main immédiate, grâce à son interface universellement connue. Quant à Zotero, il apporte la rigueur bibliographique indispensable à tout travail académique sérieux, en automatisant la gestion des références avec une précision inégalée. La recommandation principale s'oriente vers la combinaison Overleaf + Zotero, qui offre le meilleur équilibre entre qualité professionnelle, rigueur scientifique et efficacité, particulièrement adaptée aux exigences des mémoires de master et thèses.

\section{PROTOCOLE ZK-NR : RL ET POSITIONNEMENT DANS L’INVESTIGATION NUMÉRIQUE MODERNE}
\paragraph{}
Le protocole ZK-NR (Zero-Knowledge Non-Repudiation) est une architecture cryptographique modulaire en couches, axée sur la non-répudiation préservant la confidentialité pour la co-production de services numériques publics. Il combine des primitives post-quantiques (STARKs, signatures BLS à seuil, Dilithium) pour créer des preuves sécurisées, vérifiables et surtout auditables, sans jamais révéler de contenu sensible. Modélisé dans Tamarin, il s’adresse spécifiquement aux environnements réglementés (finance, e-gouvernement) en offrant des attestations juridiquement admissibles, comblant le fossé entre la sécurité cryptographique et les exigences institutionnelles.Le Trilemme CRO : Équilibre entre Confidentialité, Fiabilité et Opposabilité. La formalisation du protocole ZK-NR s’attache à son design théorique et à sa modélisation. En s’appuyant directement sur les contraintes du trilemme CRO, cette contribution détaille comment des primitives spécifiques (engagements Merkle, STARKs, etc.) sont agencées pour atteindre l’équilibre requis entre responsabilité et confidentialité. Cette phase inclut une preuve de concept et une modélisation dans l’outil Tamarin, fournissant des artefacts pratiques sans encore proposer de preuves formelles de sécurité, lesquelles sont réservées à des travaux futurs.\\
L’évolution de la cryptographie, de simple outil de protection des communications à instrument d’opposabilité juridique, transforme profondément le champ de l’investigation numérique. Désormais, les enquêteurs ne cherchent pas seulement à sécuriser ou à cacher les informations, mais à produire des preuves authentiques, vérifiables, inaltérables et légalement recevables. Les protocoles comme ZK-NR, les cadres comme CLO, ainsi que les primitives post-quantiques, ouvrent la voie à une nouvelle génération de pratiques forensiques où la preuve numérique devient incontestable devant un tribunal. Ainsi, l’investigation numérique moderne ne se limite plus à collecter des données, mais s’inscrit dans une démarche intégrée où la cryptographie devient le garant de la vérité numérique.

\section{L’UTILITÉ DE L’INVESTIGATION NUMÉRIQUE DANS LA POLICE JUDICIAIRE}
\paragraph{}
L’investigation numérique (ou digital forensic) est une discipline qui consiste à collecter, analyser, conserver et présenter des preuves numériques issues d’ordinateurs, de téléphones, de réseaux ou de tout autre support électronique, dans le but d’appuyer une enquête (judiciaire, administrative ou privée). De nos jours, elle se voit octroyée progressivement une plus grande importance dans un monde marqué par la digitalisation et la cybercriminalité et plus dans le domaine policier. L’investigation numérique permet de retrouver des traces difficiles à effacer :historiques de navigation, conversations supprimées, métadonnées, fichiers effacés mais récupérables. Elle ouvre ainsi une "scène de crime virtuelle" complémentaire à la scène physique.Les enquêtes sur le piratage informatique, les fraudes en ligne, les ransomwares, le phishing reposent directement sur ces techniques, sans investigation numérique, ces infractions resteraient impossibles à résoudre car elles laissent très peu de traces matérielles.L’analyse des adresses IP, des journaux système, des connexions réseaux permettent de remonter jusqu’au suspect. Récupération de données de géolocalisation ou de communication (SMS, WhatsApp, email) fournit des éléments d’identification et d’alibi. L’investigation permet de reconstituer une chronologie numérique : Quand un fichier a été créé, modifié, transféré ?, À quelle heure un utilisateur s’est connecté ?, Quelles données ont été effacées ou copiées ?, Ces éléments aident les enquêteurs à reconstruire le scénario d’un crime. Les procédures de collecte et de conservation (intégrité, traçabilité) assurent que les preuves numériques sont valides et utilisables devant un tribunal. Cela permet à la justice de prendre des décisions basées sur des preuves techniques solides. L’investigation numérique complète les méthodes classiques :Vidéosurveillance + analyse des communications téléphoniques. Fouilles physiques + recherche d’indices numériques. Elle donne une vision globale des faits.\\
L’investigation numérique s’est imposée comme un outil indispensable au sein de la police judiciaire, et plus particulièrement dans le contexte camerounais. Face à une criminalité moderne qui a massivement migré vers le numérique, elle est passée du statut de compétence spécialisée à celui de pilier fondamental de toute enquête criminelle. Cet atout majeur se heurte à des défis et des limites substantielles. L’explosion du volume de données, la complexité technique croissante, les contraintes juridiques et les limites matérielles et humaines, notamment la pénurie d’experts et le coût des équipements, constituent des freins réels à son efficacité maximale au Cameroun. Malgré ces obstacles, la trajectoire est tracée : il n’y a pas de retour en arrière possible. L’investigation numérique n’est plus une option, mais une nécessité pour la sécurité nationale et l’efficacité de la justice. Pour consolider ses acquis, le Cameroun doit impérativement investir dans la formation continue de ses enquêteurs, le renforcement des moyens logistiques des unités spécialisées et l’adaptation permanente de son cadre juridique.

\section{LES 10 CAS AFRICAIN LES PLUS IMPORTANT D’HACKING DURANT LES 10 DERNIERES ANNEES}
\paragraph{}
Depuis une décennie, l’Afrique s’est engagée dans une révolution numérique sans précédent, l’essor des technologies de l’information, la digitalisation des services publics et l’émergence des fintechs ont transformé les économies africaines. Mais cette ouverture numérique s’est accompagnée d’une hausse vertigineuse des cyberattaques. Selon INTERPOL (2024), le continent enregistre désormais plus de 3 000 attaques par semaine et par organisation, soit une augmentation de 300 % en dix ans. Ces attaques touchent des entreprises privées, des administrations publiques, des ONG, des universités et même des infrastructures stratégiques. Les conséquences sont multiples : pertes économiques, atteintes à la réputation, paralysie de services vitaux et fuite massive de données sensibles. Pour sélectionner ces cas, quatre critères d’évaluation permettant de mesurer la gravité et la portée de chaque attaque ont été retenu : \\
-	La taille de l’attaque : son étendue, sa durée et son niveau de complexité. \\
-	Le type d’entreprise ou d’organisation visée : publique, privée, institution financière,
ONG, etc. \\
-	Le volume de données affectées :  nature et quantité des informations compromises.\\
-	Les conséquences financières et réputationnelles pertes monétaires, sanctions ou impact social.\\
CAS: 1 - Ransomware sur Transnet (Afrique du Sud, 2021) : En juillet 2021, l'entreprise logistique publique sud-africaine Transnet a été victime d'une attaque de ransomware (rançongiciel) qui a perturbé gravement ses opérations portuaires et ferroviaires. Cette cyberattaque a entraîné l'arrêt des activités, obligeant le personnel à utiliser des méthodes manuelles et entraînant des retards considérables dans le transport des marchandises. Le groupe de pirates BlackMatter est impliqué dans l'attaque, qui a visé l'obtention d'un gain financier par l'intermédiaire de ransomware. (Taille : Attaque nationale paralysant les ports de Durban, Cape Town et Ngqura. Volume : Données logistiques et systèmes ERP chiffrés (plus de 7 To). Impact financier : Environ 60 millions USD de pertes et 3 semaines d’arrêt partiel.)\\
CAS: 2 - Breach de la CNSS (Maroc, 2025) : Breach de la CNSS" fait référence à une cyberattaque massive survenue le 8 avril 2025 contre la Caisse Nationale de Sécurité Sociale (CNSS) du Maroc, entraînant l'exfiltration de données de 2 millions de salariés et 500 000 entreprises, y comprenant des informations personnelles comme les numéros d'identification et les salaires. L'attaque a été revendiquée par un groupe se présentant comme algérien et à mise en lumière de graves vulnérabilités de sécurité, notamment l'absence de chiffrement des données. (Taille : Atteinte à une base de 2 millions de salariés et 500 000 entreprises. Volume : Données personnelles, numéros d’identification, salaires et historiques médicaux. Impact financier : Perte de confiance institutionnelle, amendes internes et coûts de remédiation élevés).\\
CAS 3: Attaque sur Eneo (Cameroun, 2024) : Enéo Cameroun a subi une cyberattaque le 29 janvier 2024, qui a perturbé ses services informatiques, notamment pour les clients prépayés et ceux utilisant des compteurs intelligents, et affecté les services de paiement en ligne. L'entreprise a communiqué sur l'incident, annonçant un retour progressif à la normale et prenant des mesures spécifiques pour les clients post payés (Taille : Perturbation des systèmes de facturation et d’achat prépayé. Volume : Données clients et journaux de transaction compromis. Impact financier : Estimé à plusieurs centaines de millions de FCFA.).\\
Et plusieurs autres cas telsque : Attaque par GhostLocker 2.0 (Egypte, 2024, Volume : Données industrielles, documents stratégiques et accès VPN volés. Impact financier : Environ 20 millions USD de rançon et pertes indirectes.) ; Le scandale Pegasus au Maroc (2020 – 2021, Volume : Données industrielles, documents stratégiques et accès VPN volés. Impact financier : Environ 20 millions USD de rançon et pertes indirectes.)etc…\\
Les dix cas présentés montrent que les attaques ne sont pas isolées, mais traduisent une transformation profonde du paysage cybercriminel. Les institutions publiques, entreprises privées et citoyens doivent désormais considérer la cybersécurité comme une responsabilité partagée. Renforcer les capacités d’investigation numérique et bâtir une souveraineté technologique ne sont plus des options : ce sont les conditions indispensables pour un développement numérique durable et sécurisé du continent africain. L’Afrique est à la croisée des chemins : son avenir numérique dépendra de sa capacité à sécuriser ses infrastructures et à former ses talents.

\end{document}